\section{Exact Lean declarations}\label{app:lean}
The following declarations use namespace
\path{CollatzPositiveProgression}.
Natural numbers include zero. The constructed starts are positive.
Here \code{oddCount k x} counts odd inputs in the first $k$ steps.
The first-hitting-time theorem requires every earlier value to
exceed one.

% Source SHA256 PositiveProgressionCoalescence.lean: 9f503a2ea7ca75c8dcdfe8b7de95618eb526b5ee5bea8c5bf69bb8c35396d58f
\noindent\begin{minipage}{\linewidth}
\paragraph{The map, its iterates, and its odd-step count.}\mbox{}\par
\noindent{\small\path{formal/PositiveProgressionCoalescence.lean}, lines 9--18.}\par
\begin{Verbatim}[formatcom=\leanappendixfont,fontsize=\scriptsize,numbers=left,firstnumber=9,numbersep=5pt,xleftmargin=18pt]
def step (x : ℕ) : ℕ :=
  if x % 2 = 0 then x / 2 else (3 * x + 1) / 2

def iterate : ℕ → ℕ → ℕ
  | 0, x => x
  | k + 1, x => iterate k (step x)

def oddCount : ℕ → ℕ → ℕ
  | 0, _ => 0
  | k + 1, x => x % 2 + oddCount k (step x)
\end{Verbatim}
\end{minipage}\par\medskip

% Source SHA256 FinitePatternCoalescence.lean: 86208b374e1df4c1ecbefcd33270853b3b0a3349fdeb0aad854ae22939729087
\noindent\begin{minipage}{\linewidth}
\paragraph{Coalescence on a common dyadic progression.}\mbox{}\par
\noindent{\small\path{formal/FinitePatternCoalescence.lean}, lines 50--53.}\par
\begin{Verbatim}[formatcom=\leanappendixfont,fontsize=\scriptsize,numbers=left,firstnumber=50,numbersep=5pt,xleftmargin=18pt]
theorem every_finite_pattern_coalesces (F : Finset ℕ) :
    ∃ K R a b : ℕ, 0 < a ∧ 0 < b ∧ ∀ n ∈ F, ∀ t : ℕ,
      iterate K (a + 2 ^ K * t + n) = b + 3 ^ R * t ∧
      oddCount K (a + 2 ^ K * t + n) = R
\end{Verbatim}
\end{minipage}\par\medskip

% Source SHA256 FinitePatternConvergence.lean: 74ccd96f004198fbd04c6a5283022cb78414fd0351f49bc4c0ad61fc1c30ffb5
\noindent\begin{minipage}{\linewidth}
\paragraph{Equal finite first hitting times for any finite pattern.}\mbox{}\par
\noindent{\small\path{formal/FinitePatternConvergence.lean}, lines 8--12.}\par
\begin{Verbatim}[formatcom=\leanappendixfont,fontsize=\scriptsize,numbers=left,firstnumber=8,numbersep=5pt,xleftmargin=18pt]
theorem every_finite_pattern_has_equal_first_hitting_times
    (F : Finset ℕ) (lower : ℕ) :
    ∃ x H R : ℕ, lower < x ∧ 1 < x ∧ ∀ n ∈ F,
      iterate H (x + n) = 1 ∧ oddCount H (x + n) = R ∧
      ∀ j < H, 1 < iterate j (x + n)
\end{Verbatim}
\end{minipage}\par\medskip

% Source SHA256 FinitePatternConvergence.lean: 74ccd96f004198fbd04c6a5283022cb78414fd0351f49bc4c0ad61fc1c30ffb5
\noindent\begin{minipage}{\linewidth}
\paragraph{Consecutive runs of every prescribed length.}\mbox{}\par
\noindent{\small\path{formal/FinitePatternConvergence.lean}, lines 32--36.}\par
\begin{Verbatim}[formatcom=\leanappendixfont,fontsize=\scriptsize,numbers=left,firstnumber=32,numbersep=5pt,xleftmargin=18pt]
theorem arbitrarily_long_consecutive_equal_first_hitting_times
    (length lower : ℕ) :
    ∃ x H R : ℕ, lower < x ∧ 1 < x ∧ ∀ n < length,
      iterate H (x + n) = 1 ∧ oddCount H (x + n) = R ∧
      ∀ j < H, 1 < iterate j (x + n)
\end{Verbatim}
\end{minipage}\par\medskip
